\chapter{第八集 \quad 天下}

\section{正文}

2400多年前的暮春，天下仍未太平，孔子却在梦中看见了久违的安宁。梦里，他和曾皙、颜回、子路这些弟子们，在春日的沂水边，身着新衣，沐浴春风，歌咏舞蹈。恍然间，对岸桃花林中有人肃立，隐约如周公。周公是孔子倾慕已久的圣人，以礼制终结了商代神权至上的治世理念，引领天下走向以人世为本的发展道路。“至哉周德”，这是目前所见，最早版本的“论语”类文献中，孔子对西周治世观的由衷推崇之语。但此刻，孔子梦中那位西周礼制开创者的身影，却已渐渐看不真切。杏林外，车马喧嚣，惊醒了孔子的梦，那是鲁哀公归城的车骑，人们纷纷传言他今日猎得麒麟。麒麟是传说中预兆圣王现世的瑞兽。

东周王室衰微，诸侯纷争不止。鲁是周公之后，理应是礼制的卫护者，如今却连鲁公也利用瑞兽来造僭越礼制的舆论之势。孔子惊觉，在有生之年里，他所怀抱的治世理想再不能如愿了。孔子理想中的天下，礼乐严明，雍容揖让。礼，是进退有止、井然有序的行为规范；乐，如舒缓优容、清和庄重的内心修养。这是曾侯乙墓编钟，由六位佩剑武士托起的巨大木架上，分三层依序悬挂65件铜钟，钟体上的铭文——标注了各钟的发音、律调、阶名，并精心标明，这些阶名与楚、周、齐、申等列国律调的对应关系。这样一套由钮钟、甬钟、镈钟组合而成的大型乐器，能演奏出由不同音色混响合成的典雅乐章，错落有致的金石之声，正暗喻着秩序整齐的礼制架构。

孔子一生起起落落，游历列国十数载，教导弟子三千，毕生正是致力于让这样雍雍穆穆的礼制思想流布天下。然而，七十高龄的他目见耳闻处，皆是礼崩乐坏。公元前548年春天的齐国，发生了一场激烈的践踏与卫护礼制之争。实权在握的齐国大夫崔杼，不顾君臣纲纪，公然杀死国君齐庄公。太史不为当权者讳，“崔杼弑其君”，短短五字，如实记录。崔杼自然不会容忍，下令杀太史、毁史册。然而太史还有两个弟弟，在那个以职为氏的年代，他们继承了史官的职责和风骨，执着地在简册上留下直笔。崔杼连杀太史一家两兄弟后，仍未能迫使史官们为其曲隐，只能作罢。不惜以性命为代价的秉笔直书，是史官们对日渐崩溃的礼制秩序的维护，也是他们在乱世中试图弥合离散的方法。

然而，在春秋三百年的动荡岁月中，旧秩序的约束力终究日渐微弱，政治格局经由血与火的洗礼，不断更迭，曾经由周天子号令的天下，转而被各怀野心的强国诸侯，搅动争霸风云。在这种情势下，宋、郑等中原小国遭兵燹之祸尤深，唯有殚精竭虑，以不同的方式辗转求生。宋国期冀通过外交斡旋来促成“弭兵”，以换取暂时的和平。宋国都城在河南商丘，今已深埋于黄河泛滥的洪积层下，经过考古勘探和发掘，我们基本了解了这段沿用千余年之久的南城墙，在东周时的模样，大部分城墙仍保留数米、甚至10米高，宽度约15米，墙体有多排纴木洞，这是修补城墙时留下的痕迹。公元前546年，这座都城曾见证了东周数百年风云中，规模最大的一场弭兵和谈。

宋都西门外，以晋、楚为首的14个诸侯国，都派出了执政的大夫参与弭战的会盟。清华简《系年》在记述此事时，虽仅有“弭天下甲兵”寥寥数字，却足以写尽沧海横流的时世里，那无数得以苟全的生灵。郑国上卿子产执政时，恰逢弭兵后的短暂和平，他从庙堂之高，俯身看到民间的芸芸众生。他深知，郑国地处中原腹心，强国环伺，朝夕有不虞之患。而子产希望邦国安宁，成为民众的依傍。因此，安内政、强军力乃当务之急。他下令“为田洫”，划田定界，防治田地不均；“作丘赋”，打破国野界限，给予乡人从军的义务和政治权利。他亦拒绝了僚属毁去乡校的建议，部分郑人对子产的咒骂传遍列国，子产自然也听说。

“苟利社稷，死生以之”。他对大夫子宽这样说道，心中的一掬赤血，终究无法剖之以示众。子产最石破天惊的一项措施是“铸刑书”，也就是把法律条文铸于铜鼎之上，公示于世人面前。列国普遍推崇礼治，仅将法制视为补充，且只有贵族才掌握律令的内容。子产却希望能让刑罚原则公开化，用成文法来规束各阶层的权利和义务，以此稳固政局。子产这一系列措施引发了部分郑国贵族的不满，也受到了列国贵族的诟责。车轮辘辘，行走在乡间的道路上，归程的子产目光悠远。夏日的木槿花娇艳而炽烈，女孩们欢声笑语，这烽火暂时消散的山水间，有人哼唱着动听的歌调，“青青子衿，悠悠我心”，“纵我不往，子宁不嗣音”。

郑地的音乐宛转悠扬，往往直抒胸臆，明快且富有生命力。虽然被正统的礼乐拥护者们，视为不入流的“淫声”，却迅速流行于各地，这不正是民间审美的抬头吗？返至家中，子产在回复叔向的责难之书中，写下了他一切作为的初心，“吾以救世也”，子产知道，战火还将再起。他肝胆俱焚，只求所有的政策能及时奏效，为郑国多争取一些时间，让那山水间平和的景象，存在的再久一些。子产被孔子称为“古之遗爱”，从子产到孔子，无不试图在这动荡的时世里，寻求重建社会秩序之法。终其一生，孔子未能见证治世理想的实践，子产也未能挽救颠坠的时局。世道艰难，但那样多的人仍在各自的位置上尽职尽责，为社稷家国死生以之，为天下大宁，匹夫怀责，这是他们的理想与信仰。

郑国国君大墓所出的一对莲鹤方壶，壶身以双龙为耳，双兽为足，腹部四角各铸一飞龙攀援欲上，通体满饰蟠螭纹，壶盖宛如一朵绽开的重瓣莲花，朵朵花瓣皆为镂空，花间站立着一只引吭振翅的鹤，昂扬向上的姿态，正如高悬于时世之上那蓬勃的朝气和理想。“鹤鸣于九皋，声闻于天”，这声声穿透混沌的清唳，在未来的历史中将会得到回应。春秋时代，小国设法求存，大国则竞逐争霸，在一众强国间，依靠实力成为天下认同的霸主。这是大国不同于小国的生存策略，黄河流域的齐国、晋国先后称霸。数十年后，齐国势头稍弱，长江流域的楚国则后来居上，与晋国形成均势。楚的崛起，与长江流域的矿产资源有关。

青铜，仍是春秋时期重要的战略资源，长江以南富集的铜、锡矿藏，原本由曾国，这一与周王室同宗的姬姓封国掌控，曾国从分封之初，除了抵御南方的淮夷之外，还承担着为周王室控制与管理铜矿资源及其运输通道的职责。曾伯漆簠的铭文中提到的“金道锡行”，指的就是连接黄河流域和长江流域的“随枣走廊”。经由崇山峻岭之间的这条狭窄山道，铜、锡等矿产资源，源源不断的北上。曾国，正是周王室这一战略生命线的安全保障。曾国国都的东南方向约200多公里处就是铜绿山，这片铜草花盛开的地方，正是东周时期规模最大的矿冶中心。已知的12个矿体中，有9个在当时便已开采。遗址中发现的矿井、炼炉、采掘和运载工具、铜锭等，展现了青铜制造业从冶炼、生产到流通的诸多环节。

长期以来，曾国都是长江与汉水之间一股强盛的力量，楚国起初完全无法与之抗衡，在随州叶家山墓地出土的斗子鼎上，郑重记载着西周初年曾国先主参与周成王会盟的经历，“丁巳，王大祓”，“己未，王赏多邦伯”，这场被曾国国君视为荣耀往事的会盟，却是楚国国君的心酸记忆。那场在岐阳举行的盛大会盟，楚国先主熊绎，也曾长途跋涉的奔赴，当他踏进仪式会场，毕恭毕敬的呈上楚地的青茅，却因实力卑弱无法正式参盟，只能与其他小国之君一起在院中看护火堆。那一夜，大殿内灯火通明，觥筹交错；冷清的庭院里，篝火熊熊，为争得列鼎而食的一席之地，南方楚人的目光里，从此留下了向北的执着。此后的楚国多代国君筚路蓝缕，方崛起于江汉西部。

春秋时期，楚国向汉水以东不断扩张，征服曾国之后，接管了铜绿山等重要矿产地，此外还掌握了丰富的锡矿资源。锡价昂贵，楚人却甚至能以锡礼器、锡箔饰随葬，获得这些重要战略资源的楚人，如虎添翼，强盛之势益发不可阻挡。数百年时移世易，楚国独霸于南方，但与姬姓的曾国联姻，对于楚国进入中原认可的秩序，仍然有着重要的意义。数代楚王都将同族的芈姓女子嫁入曾国为夫人，随州枣树林的曾国公墓中，葬有多代曾侯夫妇，191号墓中发现了夫人芈渔的器物，169号墓则出土了成套带有“随仲芈加”铭文的铜器，显示这是楚王为女儿芈加所铸的陪嫁媵器，出土的编钟上铭刻有“行相曾邦”、“余为妇为夫”等文字，记述了她曾行使国君权力的事实。

溠水汤汤，一身丧服的芈加摈退了侍从们，她是楚穆王长女曾侯宝之妻，如今丈夫早死，她与幼子面前是一个隐忧重重的曾国朝堂。芈加在楚国的支持下，决定临朝听政，从前的曾国紧紧追随周文王、周武王。随州文峰塔曾侯舆墓内出土的编钟铭文，在追忆先祖历史时，便使用了“左右文武”一语。而在附近的另一座曾侯墓出土的甬钟上，同样的语式，从属的对象却已成为了“楚王”。从“左右文武”到“左右楚王”，这一转变，正是发生在两代芈姓夫人的时代。当楚致力于对长江以南进行整合时，秦则在西部的高原绝岭间，艰难求取更大的生存空间。甘肃礼县，西汉水北岸的大堡子山上，葬着春秋早期的秦文公、静公父子。

国君大墓旁的乐器坑内，瘗埋有“秦子”镈钟、甬钟，应是文公为早死的儿子静公所铸。这套乐钟与秦公鼎、秦公簋等重器，显示出秦国对周人青铜礼器制度的忠实缵承。秦国先主原本就生活在东方，大约在西周穆王时，才迁至陇山的西边，负有为周人“保西陲”之责。由于护送周平王东迁有功，秦军得以跻身诸侯之列。继承宗周故地的秦国，被连绵山脉屏蔽于东方诸国之外，要逐鹿中原，只能通过结盟或战争的方式借道东邻的晋国。群峰寂静，大雨滂沱，刚刚自晋国凯旋的秦穆公，特地回到崤山中这处承载秦人惨痛记忆的战场。三年前，东出中原，灭滑而返的秦军在此遭遇晋军联同姜戎的伏击，数千兵士无一生还。

三年来，秦军上下无不希望一雪前恨，却屡败于晋军。直到今次，将士们渡过黄河后，将舟船尽焚，背水一战，终于大败晋军，这才敢前往崤山，为三年前阵亡的众兵士举丧。穆公于军前痛斥自己当初不听劝阻，执意发动崤之战的罪过。秦军在崤山中大哭三日，四面山谷回音阵阵，天地仿若同悲。这已是秦穆公执政的第36个年头，环顾天下，周室衰微，惟有秦与齐、楚、晋可称强国，多年来，穆公在西陲远望着东方霸主迭出，一种时不我待的急迫感，日夜催动着他。崤之战后，秦人终于认识到，东出之路已被晋国封锁，强盛之法惟有向西开拓。穆公用百里奚为左庶长，在执政的第37年，奋秦人累世经营西部之力，开地千里，称雄西戎，周天子派召公贺以金鼓。

自此，秦穆公亦被奉为春秋一霸。但崤之战也宣告着，秦晋之好已成过眼云烟，秦国东出之路暂时受阻，此后许多年，中原的政治秩序将与秦国无缘，秦国在一个相对静息的环境里，蓄力等待逐鹿中原的时机。这样的闭塞，也使得秦国的文化发展略显缓滞，列国早已抛弃的一些旧俗仍在这里延续。陕西凤翔发现的秦公一号大墓，墓主人为春秋晚期的秦景公，墓葬全长300米，是已知规模最大的先秦墓葬，墓内共发现殉葬186人，多为臣属、妃嫔以及奴婢和侍者。殉人的旧俗，铺张的墓葬，是秦国国君煊赫权力的彰示，却迥异于以周礼为纲的东方诸国。秦要参与到中原的政治格局里，还需要一场深彻的变革。

此时，以黄河岸边的函谷关为界，关城以西的秦国，蛰伏于变革的前夜之时，关城以东的局势，也正涌动起变革的浪潮，长江下游的吴国和越国，开始加入到逐鹿中原的行列。吴越壮大的基础与楚类似，也是得益于控辖的金属资源。除此外，越国更是掌握制作坚兵利器的技术。这把越王勾践剑代表了春秋时期短兵器制造的顶尖水平，该剑主要由锡青铜铸成，剑之正面，有“越王鸠浅自作用剑”八个鸟篆铭文，剑格镶有蓝色琉璃和绿松石，剑身的两面满饰黑色菱形花纹，刃薄而锋利。随着楚的强大，吴越的崛起，春秋后期的争霸主战场，逐渐从黄河流域转到长江流域，黄河流域的传统强国们则忙于应对旧有秩序瓦解的危机。其中对时局影响最深的是晋国之衰。

晋国的危机缘于公室衰微，大权旁落，六卿并立，诸卿之中，又以执政的赵氏最强。太原金胜村251号春秋晋墓，墓主赵鞅虽为正卿，却以七鼎级别随葬，俨然比肩诸侯国君。而彼时，即便是赵氏内部也人心浮动、纷争不断，卿大夫之间为维持内部团结，合力打击政敌，不得不反复的举行盟誓仪式。晋国宗庙的坛场中，司盟者手起刀落，将侍者抬来的羊杀死，殷红的牲血落入玉敦内，侍从捧着玉敦上前，参盟者逐一送出已完成书写的盟书，并以指蘸血抹于口旁。其后，侍者将杀死的祭牲埋入坎内，盟书亦被郑重地放置进去。考古工作者在盟誓遗址的400余个坎内，共发现了5000多件盟书，多以朱书写于玉圭或石圭片上，中心内容是，每个与盟人都要效忠盟主，诛讨敌对势力，并不准其重返晋地。

数量如此之多的盟书，正体现了晋国内部存在的裂隙，山盟海誓无法改变历史的走向。公元前453年，晋国被韩、赵、魏三家大夫瓜分，春秋之世历时290余年，单在《春秋》一书中所载的军事行动便有483次。春秋之初，九州天下析分为140多国，至战国之初，仅剩十余国。公室王侯、生民百姓，无不渴望结束这样的乱局。数百年乱局的一缕天光，来自一项新技术的普及。铁器，早在两周之际，炼铁技术便日趋成熟。春秋时期，秦国出土的块炼铁器数量，已远远超过其他各国。宝鸡益门村2号秦墓中，随葬有20多把异常精美的金柄剑，经鉴定，剑身的部分均为块炼铁，也就是俗称的“熟铁”。到战国时，铁器技术更是不断革新，渐成风行。相比青铜器，铁器技术具有成本低、产量大的优势，能够迅速惠及社会的各个阶层。

从考古发现的东周铁器来看，与民生密切相关的农具和手工工具最多；其次才是用于军事活动的兵器。走向普及化的铁农具和工具，极大地推动了生产力的变革，提高了人类对于自然的利用和改造能力。这样的利器，也成就了都江堰、郑国渠等各项关乎民生社稷的工程，战国时期因而迎来了社会经济的飞跃。国力与战力，固然是争夺天下的利器，但思想与文化才是深潜在世人血脉里的精髓与根本，这便是周人缔造的家国观念。现藏于中国国家博物馆的秦公簋上，刻有追述族史的铭文，“丕显朕皇祖受天命”，“宅禹蹟，十有二公在帝之坯”，自述其先祖原本生活在大禹敷布之地，凤翔秦公大墓出土石磬的铭文上，则明确提到“天子宴喜，共桓是嗣”，“高阳有灵，四方以鼏”，南方的楚人也像秦人一样，将黄帝之孙高阳追认为始祖，东方的齐国则追溯先祖为黄帝。

古史系统日益增繁，最终形成天下四方共同认可的谱系与传说，尽管干戈仍未止息，四海却已视同一家。在西方，秦与西戎之间时有摩擦，却也相互交融。位于甘肃张家川的马家塬战国墓地，正是秦文化与戎文化共生并存的写照。特殊的偏洞室墓，和随葬的陶器组合是当地戎人的传统。戎人墓主穿戴奢华，颈部戴金银半环项圈，或由各种材质珠子组成的项链，手臂上套有镶嵌繁复的金臂钏；头部饰以圆形金片和绿松石、肉红石髓珠子相间穿缀而成的发网，腰带以金银铜锡质地牌饰和珠饰组成，有的还在两侧悬挂复杂的装饰，用各类金管、肉红石髓珠、费昂斯、蜻蜓眼等串联起以高浮雕金饰或细金镶嵌玛瑙釉砂工艺的圆牌，显得斑斓绚丽。

最引人注目的，是贵族墓中随葬的大量车马，那完全按照中原样式打造的涂漆木车，车体表面，却装饰大量草原风格的金银铜铁饰件，雄鹿、猛虎和野山羊，更是属于整个欧亚草原游牧人群共有的动物形象，丰富的装饰元素充满了异族文化的张扬之美，规整连续的排列组合，又于重复中，体现华夏文化的律动与秩序。墓中随葬的铜器则体现出了秦对戎人的控制和影响力，秦与戎，这对曾经战场上的敌手，如今成了为彼此打开东西交流通途的向导。在北方，太行山东麓的中山国，由鲜虞人缔建，鲜虞，是白狄的一支，春秋时，仍被视为与中原有别的族群。在与邻近的三晋之地，发生频繁的交流、互动和摩擦之后，至战国时已渐习华风。中山王厝之冢，至今仍然高达9米，原来应是回廊环绕的三层楼阁式覆瓦建筑，顶部是享堂。

墓内随葬的青铜错金“兆域图”，绘制了陵园格局，这是迄今发现的最早的建筑设计图。该墓还出土了包括15件鼎在内的各类青铜礼器，其中一件铜方壶上錾刻着鎏金异彩的文字，言道“夫古之圣王，务在得贤，其即得民”。从墓葬建筑、随葬礼器到器物铭文，处处可见到对于周人礼制规范的遵循。而虎噬鹿器座、四龙四凤方案等青铜器，则在虎、鹿等北方式纹饰主题中，巧妙融入中原传统的龙、凤等形象，艺术与观念上的兼容并蓄，带来了生动无比的表现力。广袤大地上，处处皆是这样文化认同与文化碰撞交相辉映的动人景观。在南方，越灭吴后，成为长江下游的强国，此时，亦汇入了华夏秩序的大潮。无锡鸿山邱承墩的墓主，是地位仅次于越王的大夫，墓内随葬大量原始青瓷器，其中包括成套的编钟和编磬。

以越地的制瓷工艺，模仿青铜礼乐器型，不但象征财富与权力地位，也是越文化接触并融会华夏礼制的写照。在多元文化激荡的民族大融合时期，华夏一边坚守着文化本位，一边凝聚和吸收各族群文化中的精髓，以兼容并蓄的姿态，成为天下视野的轴心。这天下，便如一个无远弗届的同心圆，山河湖海间，处处闪耀着碰撞与交融带来的勃勃生机。

“中国”的观念，遂在这连年烽火里逐渐扩容，正如这金石乐悬是华夏礼制中最高端的“正声”，各国竞相效仿。随着时间的推移，四方与中原渐为一体，更为边远的族群，也逐渐进入时人的视野范围，“天下远近，小大若一”的理想将伴随着民族融合、国家统一的进程逐渐成为现实。

春秋战国时期，生产力的快速发展，促进了商业繁荣与城市兴盛。齐都临淄城内，曾经轻歌曼舞，百贾千工，一派繁华景象，而诸侯的争霸，又带来了对人才的广泛需求。临淄西门之外，是稷下学宫之所在，这里一度是战国时期的学术中心，曾容纳了当时诸子百家的各派学者，多达千人左右，这些学者纵横于各国之间，打破了地域的藩篱，交响出文化融会的新声。于是，新思想和新秩序的建设也随之而来。彼时的诸侯，无不反复追问，天下如何才能安定？孟子的回答是三个字：定于一。春秋战国扰攘五六百年，在中国历史上，常被当作乱世。然而，正是在这样一个时代里，中国经历了前所未有的碰撞与融合，孕育出天下大一统的基础。

进入战国时代以后，最有可能实践统一这幅蓝图的，是所谓的“七雄”，也就是函谷关以东的齐、楚、魏、燕、韩、赵六国加一个函谷关以西的秦国。此时的秦早已摈弃殉人等旧俗，但要与东边的强国抗衡，还需要更为积极的变革。秦孝公即位后，听取了魏国人商鞅的建议，在宫室盛大的栎阳城掀起了东周五百年间最风云激荡的变法。这次变法的核心，是以重典维护国家安全，进行战备积累，又以奖励军功提供上升途径，激活国力。公元前350年，秦人从栎阳迁都咸阳，商鞅在此进行了第二次变法。商鞅方升，是商鞅为统一度量衡所监制的标准量器，容积为202.15毫升，是当时的一升。这件体量不大的青铜器，却成为古中国经济秩序奠基之途上的“重器”。

商鞅两次变法，废井田，开阡陌，重农抑商，奖励军功，严革法度，推行郡县，编订户口，是一场覆盖政治、军事、经济、社会和文化制度的全方位改革，广泛调动了各种积极因素，推动秦国社会进步、经济发展、国力壮大、全面强盛。同时代的战国诸雄，虽时有一些改革作为，但却远不及秦深彻变法的决心和格局。在这风云激荡的时代里，崭新的秦国立于关中，眼望向广袤的万里河山，即将开启结束乱世，天下一统的征程。公元前328年，秦取上郡，拓地至陕北；公元前316年，秦灭巴蜀，囊括四川盆地；公元前272年，秦灭义渠，攻克了黄土高原最大的戎狄势力。秦国就这样，以关中盆地为营，分别向北、西、南三个方向扩展战略纵深，随后剑锋东指，匡合天下。

在商鞅变法的百年之后，秦的统一大业已势如破竹。湖北云梦，本为楚地，被秦人占据后，成为攻灭楚国的战略要地。埋葬在睡虎地的喜，是秦军中的一位下级军吏，亲身经历了秦完成统一的最后阶段。随葬的《编年记》载录着他的这段履历，嬴政在位的第十三年，喜从军，十五年从平阳军，十七年，攻韩，十八年，攻赵，十九年，南郡备警，廿一年，韩王死，廿二年，攻魏梁，廿三年，攻荆，楚昌文君死。甘肃礼县，海拔1867米的四角子山之巅，以中央方形土台，及其上半地穴房址为中心，四周对称分布着平行的双排建筑。这是秦始皇统一六国后，在秦人祖居之地“犬丘”修建的祭天场所。公元前221年，秦始皇熔化天下兵器，铸十二铜人，颁发诏版，宣告天下“廿六年，皇帝尽并兼天下诸侯，黔首大安”。三晋的强兵劲弩，楚的坚守抵抗，齐的合纵连横，燕的舍生忘死，都是向着天下一统驱驰的时代侧影。

秦，终结了数百年的纷争，延续和实现了周人的家国理想，将华夏九州推入统一多民族国家的时代。站在江河入海的时代，回顾奔涌的历史长流，是蹈海踏浪而来的文明记忆。曾有人在黑夜里，擎起天地间的第一簇火种，曾有人在漫天星斗下，隔着山川湖海遥望。人聚成邑，邑聚为国，曾有古国在大江大河间初铸辉煌，自九州之中，酿就四方众望、人心所向的风潮。经夏商周而立秦汉，生长在世界东方的古代中国文明，跋涉数千年岁月而来，穿梭在生死枯荣的古文明之林间，历风尘霜雪而初心未改。蹈蹚出一条不同于世界其他文明体的发展道路，创造出灿烂盛大的中华文明。《何以中国》从始至终，亦是站在考古学搭起的阶梯上，向千百代来成就今日之中国的祖先们致敬！自先人起，便未更改对于这片土地的深情，以及对于和平与统一的刻骨记忆，让我们在这片祖先们曾守望过的天下，携手同心，去往那光明盛放的前方！

